% GENERATED FILE. Edit canonical lessons, then run npm run generate:books.
% canonical-source-hash: fnv1a64:bfc3928b70a7b92c
% canonical-lessons: TA-C215-pencil, TA-C215-akarati, TA-C215-kattanam, TA-C215-katitam, TA-C215-urai

\chapter{Pencil, Dictionary, Fee, Letter, Envelope}
\label{ch:ta-pencil-dictionary-fee-letter-envelope}
\hlchaptermodality{215}

\begin{chapteropening}
\textbf{By the end of this chapter:} \emph{I can say a pencil, a dictionary, a fee, a letter and an envelope.}

Complete the last lesson of chapter 215: I can say a pencil, a dictionary, a fee, a letter and an envelope.
\end{chapteropening}

\section[\texorpdfstring{\ta{பென்சில்} (peṉcil)}{பென்சில் (peṉcil)}]{\ta{பென்சில்} (peṉcil) — a pencil}

\label{lesson:TA-C215-pencil}



\begin{quote}
Before the new one: say the Tamil for flour; batter, then the Tamil for bread.
\end{quote}

\subsection*{You'll want to know: \ta{பென்சில்}}
\textbf{\ta{பென்சில்}} — \emph{peṉcil} — \textquotedblleft{}a pencil\textquotedblright{}.

\begin{grammarlens}[title={What you've built}]
One more word for this chapter.
\end{grammarlens}

\subsection*{Your turn}
\begin{itemize}
\raggedright
  \item \emph{Say it:} \emph{peṉcil}
  \item \emph{Say it:} \emph{peṉcil}, once more
  \item \emph{Say it:} say \emph{peṉcil}
  \item \emph{Recall:} say the Tamil for flour; batter, then the Tamil for bread, then say \emph{peṉcil} again
\end{itemize}

\begin{tcolorbox}[breakable,colback=teal!4,colframe=teal!35!black,title={Before you move on}]
What does \ta{பென்சில்} mean? (\textquotedblleft{}A pencil\textquotedblright{}.) Say it once more.
\end{tcolorbox}
\section[\texorpdfstring{\ta{அகராதி} (akarāti)}{அகராதி (akarāti)}]{\ta{அகராதி} (akarāti) — a dictionary}

\label{lesson:TA-C215-akarati}



\begin{quote}
Before the new one: say the Tamil for bread, then the Tamil for a pencil.
\end{quote}

\subsection*{You'll want to know: \ta{அகராதி}}
\textbf{\ta{அகராதி}} — \emph{akarāti} — \textquotedblleft{}a dictionary\textquotedblright{}.

\begin{grammarlens}[title={What you've built}]
One more word for this chapter.
\end{grammarlens}

\subsection*{Your turn}
\begin{itemize}
\raggedright
  \item \emph{Say it:} \emph{akarāti}
  \item \emph{Say it:} \emph{akarāti}, once more
  \item \emph{Say it:} say \emph{akarāti}
  \item \emph{Recall:} say the Tamil for bread, then the Tamil for a pencil, then say \emph{akarāti} again
\end{itemize}

\begin{tcolorbox}[breakable,colback=teal!4,colframe=teal!35!black,title={Before you move on}]
What does \ta{அகராதி} mean? (\textquotedblleft{}A dictionary\textquotedblright{}.) Say it once more.
\end{tcolorbox}
\section[\texorpdfstring{\ta{கட்டணம்} (kaṭṭaṇam)}{கட்டணம் (kaṭṭaṇam)}]{\ta{கட்டணம்} (kaṭṭaṇam) — a fee, a charge}

\label{lesson:TA-C215-kattanam}



\begin{quote}
Before the new one: say the Tamil for a pencil, then the Tamil for a dictionary.
\end{quote}

\subsection*{You'll want to know: \ta{கட்டணம்}}
\textbf{\ta{கட்டணம்}} — \emph{kaṭṭaṇam} — \textquotedblleft{}a fee, a charge\textquotedblright{}.

\begin{grammarlens}[title={What you've built}]
One more word for this chapter.
\end{grammarlens}

\subsection*{Your turn}
\begin{itemize}
\raggedright
  \item \emph{Say it:} \emph{kaṭṭaṇam}
  \item \emph{Say it:} \emph{kaṭṭaṇam}, once more
  \item \emph{Say it:} say \emph{kaṭṭaṇam}
  \item \emph{Recall:} say the Tamil for a pencil, then the Tamil for a dictionary, then say \emph{kaṭṭaṇam} again
\end{itemize}

\begin{tcolorbox}[breakable,colback=teal!4,colframe=teal!35!black,title={Before you move on}]
What does \ta{கட்டணம்} mean? (\textquotedblleft{}A fee, a charge\textquotedblright{}.) Say it once more.
\end{tcolorbox}
\section[\texorpdfstring{\ta{கடிதம்} (kaṭitam)}{கடிதம் (kaṭitam)}]{\ta{கடிதம்} (kaṭitam) — a letter (to post)}

\label{lesson:TA-C215-katitam}



\begin{quote}
Before the new one: say the Tamil for a dictionary, then the Tamil for a fee.
\end{quote}

\subsection*{You'll want to know: \ta{கடிதம்}}
\textbf{\ta{கடிதம்}} — \emph{kaṭitam} — \textquotedblleft{}a letter (to post)\textquotedblright{}.

\begin{grammarlens}[title={What you've built}]
One more word for this chapter.
\end{grammarlens}

\subsection*{Your turn}
\begin{itemize}
\raggedright
  \item \emph{Say it:} \emph{kaṭitam}
  \item \emph{Say it:} \emph{kaṭitam}, once more
  \item \emph{Say it:} say \emph{kaṭitam}
  \item \emph{Recall:} say the Tamil for a dictionary, then the Tamil for a fee, then say \emph{kaṭitam} again
\end{itemize}

\begin{tcolorbox}[breakable,colback=teal!4,colframe=teal!35!black,title={Before you move on}]
What does \ta{கடிதம்} mean? (\textquotedblleft{}A letter (to post)\textquotedblright{}.) Say it once more.
\end{tcolorbox}
\section[\texorpdfstring{\ta{உறை} (uṟai)}{உறை (uṟai)}]{\ta{உறை} (uṟai) — an envelope; a cover}

\label{lesson:TA-C215-urai}



\begin{quote}
Before the new one: say the Tamil for a fee, then the Tamil for a letter.
\end{quote}

\subsection*{You'll want to know: \ta{உறை}}
\textbf{\ta{உறை}} — \emph{uṟai} — \textquotedblleft{}an envelope; a cover\textquotedblright{}.

\begin{grammarlens}[title={What you've built}]
One more word for this chapter.
\end{grammarlens}

\subsection*{Your turn}
\begin{itemize}
\raggedright
  \item \emph{Say it:} \emph{uṟai}
  \item \emph{Say it:} \emph{uṟai}, once more
  \item \emph{Say it:} say \emph{uṟai}
  \item \emph{Recall:} say the Tamil for a fee, then the Tamil for a letter, then say \emph{uṟai} again
\end{itemize}

\begin{tcolorbox}[breakable,colback=teal!4,colframe=teal!35!black,title={Before you move on}]
What does \ta{உறை} mean? (\textquotedblleft{}An envelope; a cover\textquotedblright{}.) Say it once more.
\end{tcolorbox}
